% Part: first-order-logic
% Chapter: natural-deduction
% Section: derivations

\documentclass[../../../include/open-logic-section]{subfiles}

\begin{document}

\iftag{FOL}
      {\olfileid{fol}{ntd}{der}}
      {\olfileid{pl}{ntd}{der}}

\olsection{\usetoken{P}{derivation}}

\begin{explain}
অনুমিতি কী তা বলেছি এবং অনুমান-বিধিগুলিও দিয়েছি। এগুলি থেকে
স্বাভাবিক নিষ্পাদনের !!^{derivation}s আরোহীভাবে গঠিত হয়: প্রত্যেক
!!{derivation} হয় একাই একটি অনুমিতি, নয়তো এক, দুই বা তিনটি
!!{derivation}s-এর পরে একটি সঠিক অনুমান দিয়ে গঠিত।
\end{explain}

\begin{defn}[!!^{derivation}]
অনুমিতি~$\Gamma$ থেকে কোনো !!a{sentence}~$!A$-এর একটি
\emph{!!{derivation}} হলো !!{sentence}s-এর এমন একটি সসীম বৃক্ষ যা
নিচের শর্তগুলি পূরণ করে:
\begin{enumerate}
\item বৃক্ষের সর্বোচ্চ !!{sentence}s হয়~$\Gamma$-তে আছে, নয়তো
  বৃক্ষের কোনো অনুমান সেগুলিকে !!{discharged} করে।
\item বৃক্ষের সর্বনিম্ন !!{sentence} হলো~$!A$।
\item বৃক্ষের নিচে থাকা বাক্য~$!A$ ছাড়া প্রত্যেক !!{sentence} কোনো
  অনুমান-বিধির সঠিক প্রয়োগের পূর্বধারণা; সেই বিধির সিদ্ধান্তটি ওই
  !!{sentence}-এর ঠিক নিচে বৃক্ষে থাকে।
\end{enumerate}
তখন $!A$-কে !!{derivation}-টির \emph{সিদ্ধান্ত} এবং $\Gamma$-কে তার
!!{undischarged} অনুমিতি বলি।

$\Gamma$ থেকে~$!A$-এর কোনো !!a{derivation} থাকলে বলি, $!A$ হলো
$\Gamma$ থেকে \emph{!!{derivable}}; প্রতীকে: $\Gamma \Proves
!A$। $!A$-এর এমন কোনো !!a{derivation} থাকলে যার প্রত্যেক অনুমিতি
!!{discharged}, লিখি~$\Proves !A$।
\end{defn}

\begin{ex}
প্রত্যেক অনুমিতি নিজেই !!a{derivation}। যেমন, $!A$ নিজে একটি
!!a{derivation}, আর $!B$ নিজেও তাই। এগুলিতে, ধরা যাক,
$\Intro{\land}$ বিধি প্রয়োগ করে একটি নতুন !!{derivation} পাই:
\begin{prooftree}
\AxiomC{$!A$}
\AxiomC{$!B$}
\RightLabel{\Intro{\land}}
\BinaryInfC{$!A \land !B$}
\end{prooftree}
এই বিধিগুলি সাধারণ হওয়ার কথা: এর~$!A$ ও~$!B$-এর জায়গায় যেকোনো
!!{sentence}s, যেমন $!C$ ও~$!D$, বসাতে পারি। তখন সিদ্ধান্ত হবে
$!C \land !D$; তাই
\begin{prooftree}
\AxiomC{$!C$}
\AxiomC{$!D$}
\RightLabel{\Intro{\land}}
\BinaryInfC{$!C \land !D$}
\end{prooftree}
একটি সঠিক !!{derivation}। অবশ্যই অনুমিতিগুলির স্থানও বদলাতে পারি, যেন
$!D$~$!A$-এর এবং $!C$~$!B$-এর ভূমিকা পালন করে। সুতরাং
\begin{prooftree}
\AxiomC{$!D$}
\AxiomC{$!C$}
\RightLabel{\Intro{\land}}
\BinaryInfC{$!D \land !C$}
\end{prooftree}
এটিও একটি সঠিক !!{derivation}।

এবার আর-একটি বিধি, ধরা যাক $\Intro{\lif}$, প্রয়োগ করতে পারি। এটি
একটি শর্তাধীন বাক্যে পৌঁছতে এবং সেই শর্তাধীন বাক্যের পূর্বপদের সঙ্গে
অভিন্ন যেকোনো অনুমিতিকে !!{discharge} করতে দেয়। তাই নিচের দুটিই
সঠিক !!{derivation}s:
\begin{prooftree}
\AxiomC{$\Discharge{!C}{1}$}
\AxiomC{$!D$}
\RightLabel{\Intro{\land}}
\BinaryInfC{$!C \land !D$}
\DischargeRule{\Intro{\lif}}{1}
\UnaryInfC{$!C \lif (!C \land !D)$}
\DisplayProof\bottomAlignProof
\AxiomC{$!C$}
\AxiomC{$\Discharge{!D}{1}$}
\RightLabel{\Intro{\land}}
\BinaryInfC{$!C \land !D$}
\DischargeRule{\Intro{\lif}}{1}
\UnaryInfC{$!D \lif (!C \land !D)$}
\end{prooftree}
এগুলি যথাক্রমে দেখায় যে $!D \Proves !C \lif (!C \land !D)$ এবং
$!C \Proves !D \lif (!C \land !D)$।

মনে রেখো, অনুমিতি অবমুক্ত করা একটি অনুমতি, বাধ্যবাধকতা নয়: অনুমিতি
অবমুক্ত করতেই হবে এমন নয়। বিশেষত, অনুমিতিগুলি !!{derivation}-এ না
থাকলেও বিধি প্রয়োগ করা যায়। যেমন, নিচেরটি বৈধ, যদিও !!{discharged}
করার মতো কোনো অনুমিতি~$!A$ নেই:
\begin{prooftree}
  \AxiomC{$!B$}
  \DischargeRule{\Intro{\lif}}{1}
  \UnaryInfC{$!A \lif !B$}
\end{prooftree}
\end{ex}

\end{document}
